\documentclass[a4paper,addpoints]{article}
\usepackage[utf8]{inputenc}
\usepackage{amsfonts,amssymb,amsmath, amsthm}
\usepackage{xcolor}
\usepackage{pgf,tikz,pgfplots}
\pgfplotsset{compat=1.15}
\usepgfplotslibrary{fillbetween}
\usepackage{mathrsfs}
\usetikzlibrary{arrows}
\usetikzlibrary{calc}
\usepackage{tabu}

\usepackage{tikz}
\usetikzlibrary {shapes , arrows , fit , calc, positioning}
\tikzset{box/.style={draw, rectangle, rounded corners, thick, node distance =7 em, text width =6 em, text centered, minimum height =2.5 em}}
\tikzset{line/.style={draw, thick, -latex'}}

\usepackage{tikz-uml}
\tikzumlset{actor below=0.4cm, fill usecase=white, fill class=white, text=black, draw=black}

\begin{document}

In Lostly ogni \emph{utente} è identificato da nome, email e password e può pubblicare più \emph{segnalazioni}. Una segnalazione ha una descrizione, un luogo, una data, una foto e uno stato, e appartiene a una \emph{categoria}; può essere di due tipi: oggetto \emph{smarrito} oppure oggetto \emph{ritrovato}. Un oggetto smarrito può essere messo in relazione con più oggetti ritrovati (e viceversa) tramite una \emph{corrispondenza}, che ha una data di segnalazione e uno stato. Ogni corrispondenza genera delle \emph{notifiche} destinate agli utenti coinvolti, e gli utenti possono scambiarsi dei \emph{messaggi} relativi a una corrispondenza per organizzare la restituzione. L'\emph{amministratore} è un utente che può gestire utenti e segnalazioni.

\vspace{1cm}
\begin{center}
\resizebox{\textwidth}{!}{%
\begin{tikzpicture}

\umlclass[x=-2]{Utente}{Nome : string \\ Email : string \\ Password : string}{registrati() : void \\ accedi() : void}
\umlclass[x=-10]{Amministratore}{}{gestisciUtenti() : void \\ gestisciSegnalazioni() : void}
\umlinherit{Amministratore}{Utente}

\umlclass[x=6]{Segnalazione}{Descrizione : string \\ Luogo : string \\ Data : date \\ Foto : string \\ Stato : string}{modifica() : void \\ aggiornaStato(s : string) : void}
\umlassoc[mult1=1, mult2=*, pos=0.8]{Utente}{Segnalazione}

\umlclass[x=14]{Categoria}{Nome : string}{}
\umlassoc[mult1=*, mult2=1]{Segnalazione}{Categoria}

\umlsimpleclass[x=1, y=-6]{Smarrito}
\umlsimpleclass[x=9, y=-4.5]{Ritrovato}
\umlinherit{Smarrito}{Segnalazione}
\umlinherit{Ritrovato}{Segnalazione}

\umlHVassoc[name=assoc, mult1=*, mult2=*, pos=1.7]{Smarrito}{Ritrovato}
\umlassocclass[x=5.3, y=-9]{Corrispondenza}{assoc-1}{DataSegnalazione : date \\ Stato : string}{}

\umlclass[x=-9, y=-6]{Notifica}{Testo : string \\ Data : date \\ Letta : bool}{}
\umlassoc[mult1=1, mult2=*]{Utente}{Notifica}
\umlassoc[mult1=*, mult2=1]{Notifica}{Corrispondenza}

\umlclass[x=-5, y=-10.5]{Messaggio}{Testo : string \\ Data : date}{}
\umlassoc[mult1=1, mult2=*]{Utente}{Messaggio}
\umlassoc[mult1=*, mult2=1]{Messaggio}{Corrispondenza}

\end{tikzpicture}}
\end{center}

\end{document}
